\section{Maximal-intersection ownership aggregation}
\label{sec:affine-ownership-maximal-intersection}

We continue the affine catalogue notation of
Section~\ref{sec:affine-inverse-period-catalogue}.  This section proves an
unconditional compression theorem for the non-arm catalogue.  It does not
prove the \(abc\) conjecture: the final strict comparison isolated below
remains open.

Let \(X\) be the selected canonical parameter points in a square of side
difference \(N\).  Write \(\kappa(x)=(k_U(x),k_V(x),k_W(x))\) for the
powerful-kernel triple of \(x\), and order triples by coordinatewise
divisibility.  For distinct \(x,y\in X\), set
\[
 \gamma(x,y)=\gcd_{\rm coord}(\kappa(x),\kappa(y)).
\]
Let \(\Gamma\) be the distinct values \(\gamma(x,y)\) having coordinate
product greater than \(N^2\) and a non-arm supporting line, and let
\(\mathcal M\) be the divisibility-maximal elements of \(\Gamma\).  Put
\[
 \operatorname{Supp}(\mu)=\{x\in X:\mu\preceq\kappa(x)\},\qquad
 r_\mu=|\operatorname{Supp}(\mu)|.
\]

\begin{proposition}[Maximal supports and the reduced skeleton]
\label{prop:affine-maximal-support-skeleton}
Every repeated non-arm large label divides a maximal top on the same unique
non-arm line.  For \(\mu\in\mathcal M\), one has \(r_\mu\ge2\), and every
pair of distinct points in \(\operatorname{Supp}(\mu)\) has exact top
\(\mu\).  Consequently distinct maximal supports meet in at most one point
and
\[
 \sum_{\mu\in\mathcal M}\binom{r_\mu}{2}\le\binom{|X|}{2}.
 \tag{\ref{prop:affine-maximal-support-skeleton}.1}
\]
Moreover, the reduced class skeleton consisting of one loop for every
kernel class of multiplicity at least two and one edge for every pair of
singleton classes is cofinal in \(\Gamma\), and hence has exactly the same
maximal top values.
\end{proposition}

\begin{proof}
A repeated label is contained in the exact gcd top of any two of its
support points.  Finiteness supplies a maximal top above it.  The top is
itself a large label, so large-label line uniqueness keeps every comparable
top on the original non-arm line.  If two points support a maximal top
\(\mu\), their exact top lies above \(\mu\) and on that line; maximality
forces equality.  Two common points of distinct maximal supports would
therefore have two different exact tops, proving the intersection and pair
count assertions.

For cofinality, consider an exact pair top omitted by the reduced skeleton.
At least one endpoint belongs to a class \(\kappa\) of multiplicity at least
two.  Its retained loop has top \(\kappa\), which dominates the omitted gcd.
Every point in this class supports the omitted large label, so the loop is
on the same unique non-arm line and is eligible.  A cofinal subset of a
finite poset has the same maximal elements: dominate a whole-poset maximum
by a subset element, and dominate any whole-poset extension of a subset
maximum once more inside the subset; antisymmetry then gives equality.
\end{proof}

Fix a total order on \(\mathcal M\), and assign every repeated non-arm label
\(\lambda\) to the first maximal top \(o(\lambda)\) divisible by it.  Define
\[
 \mathcal O_\mu=\{\lambda:o(\lambda)=\mu\},\qquad
 S_\mu=\sum_{\lambda\in\mathcal O_\mu}\frac{w_\lambda}{T_\lambda^2},\qquad
 E_\mu=\sum_{\lambda\in\mathcal O_\mu}w_\lambda(n_\lambda-1)^3.
\]
These are exact partitions of \(S_{\rm non}\) and \(E_{\rm non}\).  If
\(\lambda\in\mathcal O_\mu\), then
\[
 n_\lambda\ge r_\mu,\qquad T_\lambda\mid T_\mu.
 \tag{\ref{prop:affine-maximal-support-skeleton}.2}
\]
The second relation is in this direction: for \(d\mid m\),
\(d/\gcd(d,a)\mid m/\gcd(m,a)\) coordinatewise.  It does not permit
\(T_\lambda^{-2}\) to be bounded above by \(T_\mu^{-2}\).

On the line of \(\mu\), let
\[
 Q(\mu)=\sum_{\substack{d\preceq\mu\\D_d>N^2}}\frac{w_d}{T_d^2},\qquad
 B_\mu\ge S_\mu,\qquad
 \mathcal H_3=\sum_{\mu\in\mathcal M}\frac{B_\mu}{(r_\mu-1)^3}.
\]
One valid owner-independent choice has
\[
 B_\mu=\min\left\{F(\mu,A_\mu),\frac{C_\mu^2G_\mu}{N^4},
              \frac{C_\mu^2H(G_\mu)}{N^2}\right\},\qquad
 S_\mu\le Q(\mu)\le B_\mu,
 \tag{\ref{prop:affine-maximal-support-skeleton}.3}
\]
with the notation of Theorem~\ref{thm:affine-hybrid-pair-catalogue}.

\begin{theorem}[Ownership Cauchy inequality and strict ray closure]
\label{thm:affine-ownership-cauchy}
One has
\[
 S_{\rm non}^2\le E_{\rm non}\mathcal H_3.
 \tag{\ref{thm:affine-ownership-cauchy}.1}
\]
Writing \(K=(B+1)(C+1)\), the nonempty non-arm domain satisfies the strict
bounds
\[
 S_{\rm non}<KN\mathcal H_3,\qquad
 E_{\rm non}<(KN)^2\mathcal H_3.
 \tag{\ref{thm:affine-ownership-cauchy}.2}
\]
For an empty domain the corresponding weak inequalities hold with both
sides zero.
\end{theorem}

\begin{proof}
For each maximal top, catalogue coverage and support containment give
\[
 S_\mu\le B_\mu,\qquad
 S_\mu\le\sum_{\lambda\in\mathcal O_\mu}w_\lambda
       \le\frac{E_\mu}{(r_\mu-1)^3}.
\]
Multiplication and Cauchy--Schwarz over the finite index set
\(\mathcal M\), rather than over labels or point pairs, prove the first
claim.  The signed-ray theorem gives termwise
\((n_\lambda-1)^3T_\lambda^2<KN\).  Positivity of
\(w_\lambda/T_\lambda^2\) makes the summed inequality
\(E_{\rm non}<KN S_{\rm non}\) strict whenever the exact summation domain is
nonempty.  In that case \(S_{\rm non}>0\), and the owner cap of the nonempty
owner set also gives \(\mathcal H_3>0\).  Combining the two inequalities and
dividing by \(S_{\rm non}\) proves the strict conclusions.
\end{proof}

\begin{proposition}[Full-catalogue multiplicity is paid by energy]
\label{prop:affine-catalogue-multiplicity-energy}
The maximal-top catalogues satisfy
\[
 \sum_{\mu\in\mathcal M}Q(\mu)\le E_{\rm non}.
 \tag{\ref{prop:affine-catalogue-multiplicity-energy}.1}
\]
Thus the exact choice \(B_\mu=Q(\mu)\) gives
\(\mathcal H_3\le\sum_\mu Q(\mu)\le E_{\rm non}\).
\end{proposition}

\begin{proof}
For a repeated non-arm label \(\lambda\), let
\(m_\lambda=\#\{\mu\in\mathcal M:\lambda\preceq\mu\}\).  Every divisor term
of \(Q(\mu)\) is an actual repeated label on the same unique line, so finite
summation gives the exact identity
\[
 \sum_{\mu\in\mathcal M}Q(\mu)=
 \sum_{\lambda\ {\rm repeated,nonarm}}
       \frac{w_\lambda}{T_\lambda^2}m_\lambda.
\]
Choose one realizing pair for every top counted by \(m_\lambda\).  These
pairs lie in the \(n_\lambda\)-point support and are distinct, because one
pair has one exact gcd top.  Hence
\(m_\lambda\le\binom{n_\lambda}{2}\).  For \(n_\lambda\ge2\) and
\(T_\lambda\ge1\),
\[
 \frac1{T_\lambda^2}\binom{n_\lambda}{2}\le(n_\lambda-1)^3.
\]
Multiplying by \(w_\lambda\) and summing proves the claim.  Finally
\(r_\mu\ge2\), so the cubic support denominator only decreases \(Q(\mu)\).
\end{proof}

Let \(\mathcal A_{\rm arm}\) denote the three proved arm caps and let
\(\widehat Q_{\rm dir}\) be the sum of the pairwise-coprime direction-lcm
envelopes.  Exact label containment in those envelopes and the preceding
theorem give
\[
 E_{\rm sh}\le
 \min\{KN\widehat Q_{\rm dir},(KN)^2\mathcal H_3\}
 +\mathcal A_{\rm arm}.
 \tag{\ref{thm:affine-ownership-cauchy}.3}
\]
Since \(J=A_1+\Omega\), \(W=A_0-\Omega\), and
\(J^3\le W^2E_{\rm sh}\), every selected configuration with \(W>0\) obeys
\[
 \left(\frac{A_1+\Omega}{A_0-\Omega}\right)^3
 \le\min\left\{KN\frac{\widehat Q_{\rm dir}}{A_0-\Omega},
  (KN)^2\frac{\mathcal H_3}{A_0-\Omega}\right\}
 +\frac{\mathcal A_{\rm arm}}{A_0-\Omega}.
 \tag{\ref{thm:affine-ownership-cauchy}.4}
\]
This is the ownership-preserving normalized gate.

For completeness, in the canonical consecutive-coefficient setting
\(1\le B\), \(C=B+1\), the arithmetic divisor enumeration gives, for every
\(\varepsilon>0\),
\[
 \begin{aligned}
 |\mathcal M|&\ll_\varepsilon(CN)^\varepsilon
        \min\{N^2,CN^{11/6}\},\\
 \sum_{\mu\in\mathcal M}B_\mu&\ll_\varepsilon
 K(CN)^\varepsilon\min\{N^5,CN^{29/6}\},
 \end{aligned}
 \tag{\ref{thm:affine-ownership-cauchy}.5}
\]
provided \(S_\mu\le B_\mu\le F(\mu,A_\mu)\).  On
\(r_\mu-1\ge R_0\ge1\), the pressure has the further factor \(R_0^{-4}\).
Indeed a closest support pair has step \(q_\mu\) satisfying
\(q_\mu(r_\mu-1)L\le N\); the fixed-direction candidates divide the three
integers \(q_\mu|A_Z|\); and the original, before \(R\)-free reduction,
coefficient is finite-to-one on each side of the direction square.  Here the
three direction coefficients are pairwise coprime after removing the primes
of \(\operatorname{rad}(BC)\): their pairwise gcds divide \(B\), \(C\), and
\(C-B=1\).  The square-divisor excess estimate then gives
\(O(CL^{5/6})\) ordinary
directions on shell \(L\), and summation gives the displayed exponents.

\paragraph{Certified boundaries.}
The following statements are refuted only with the premises stated here.
All examples are exact finite certificates.
\begin{enumerate}
\item The implication \(T_\lambda\ge T_\mu\) for every owned label is false
under canonical admissibility, non-arm maximal ownership, and full catalogue
containment.  At \(B=7,C=8,M=400,N=399\), the unique top
\(\mu=(9,121,529)\) has \(T_\mu=3\), while its proper owned label
\((3,121,529)\) has period one; the exact catalogue has two terms and
\(Q(\mu)=445280/3\).
\item The assertion that one kernel class belongs to at most one maximal top
is false under canonical admissibility and global divisibility maximality.
At \(B=5,C=6,M=254,N=253\), the singleton class
\((1,841,43681)\) lies in the two incomparable maximal supports with tops
\((1,841,121)\) and \((1,841,361)\), whose intersection is exactly that one
point.
\item Raw pair enumeration is not ownership preserving.  At
\(B=4,C=5,M=390,N=389\), all three pairs of three canonical points have the
same top \((49,361,9)\).  Its support has size three, its period-one
catalogue has mass \(86184\), and equality holds in the local cubic Cauchy
bound after the single top is deduplicated.
\item Neither \(Q(\mu)\le w_\mu\) nor \(Q(\mu)\le2w_\mu\) follows from
canonical admissibility, non-arm maximal ownership, and full catalogue
coverage.  At \(B=55123,C=55124,M=3,N=2\), the two-point top
\(\mu=(1,1,55125)\) has 34 period-one terms,
\[
 Q(\mu)=55122,\qquad w_\mu=25200,\qquad
 \frac{Q(\mu)}{w_\mu}=\frac{9187}{4200}>2.
\]
The explicit three-term cap equals \(55125\), with inflation \(35/16>2\).
\item Maximality, ownership, cap coverage, \(r_\mu\ge2\), and linear support
alone imply neither a strict Cauchy saving nor a linear number of tops.  The
abstract four-vertex construction assigns an edge-specific prime square to
each pair, realizes all six pair tops, and has
\(S_{\rm non}=E_{\rm non}=\mathcal H_3\).  This item is an abstract ledger
certificate and does not claim canonical affine realizability.
\end{enumerate}
The first four examples are canonical admissible local selections, but they
are not asserted to satisfy the exceptional-point inequality.  They refute
the listed local implications and no stronger statement that adds
exceptional-point geometry.  The fifth omits the affine divisibility and
powerful-excess filters.  None of these certificates closes the affine
route.

\paragraph{Open gate.}
The unresolved task is to prove the strict reverse of
\[
 (A_1+\Omega)^3\le(A_0-\Omega)^2
 \left[\min\{KN\widehat Q_{\rm dir},(KN)^2\mathcal H_3\}
       +\mathcal A_{\rm arm}\right]
 \tag{\ref{thm:affine-ownership-cauchy}.6}
\]
for every hypothetical exceptional selection, or to obtain an actual
standard \(abc\) counterexample.  This requires a canonical low-support
supersaturation or catalogue-sparsity theorem coupling the normalized
incidence, \(\mathcal H_3\), and the arm term.  The finite scan found no
same-line or same-direction multiplicity in its tested range, but these are
explicitly finite no-hit observations and are not used as theorems.

\paragraph{Formalization boundary.}
The companion Lean module proves 24 finite or algebraic theorems and prints
the axiom dependencies of all 24.  It formalizes finite-poset cofinality,
abstract maximal-support and codegree arguments, owner sums, Cauchy
aggregation, strict-ray closure, the polynomial normalized gate, abstract
tree subtraction, and the numerical cores of the boundary certificates.
It does not formalize the affine realization of the maximal cover, the
general period-divisibility lemma, the exact catalogue double count and its
support-pair injection, the direction Euler products and tails, the
spanning-tree realization, the closest-pair shell counts, the arm and
weighted H\"older instantiations, or the complete canonical-kernel
verification of the finite examples.  The analytic bounds above remain
paper proofs; the complete finite canonical premises are independently
replayed by the accompanying Python certificate.
